% Standalone export: контекстная диаграмма IDEF0 (A0) — постановка задачи
\documentclass[tikz,border=8pt]{standalone}
\usepackage[utf8]{inputenc}
\usepackage[T2A]{fontenc}
\usepackage[russian]{babel}
\usetikzlibrary{arrows.meta, positioning, calc}

\begin{document}
\begin{tikzpicture}[
  font=\normalsize,
  arr/.style={-{Stealth[length=2.8mm]}, thick},
  lbl/.style={font=\small, inner sep=1pt},
  arrlen/.store in=\arrlen,
  arrlen=12mm,
  outlen/.store in=\outlen,
  outlen=12mm,
  stub/.store in=\stub,
  stub=14mm,
]
% Центральный блок (функция)
\node[draw, very thick, minimum width=10.5cm, minimum height=4.2cm, align=center]
  (sys) {Хранение и анализ\\ статистики игроков\\ и команд NBA};
\node[anchor=south east, inner sep=3pt, font=\small] at (sys.south east) {A0};

% --- ВХОДЫ (слева) ---
\foreach \y/\t in {
    17/{Первичная статистика\\ матчей и игроков},
    0/{Справочные данные\\ (игроки, команды, сезоны)},
    -17/{Запросы аналитиков\\ и пользователей}} {
  \node[lbl, anchor=east, align=right, text width=36mm]
    at ([xshift=-\arrlen,yshift=\y mm]sys.west) {\t};
  \draw[arr] ([xshift=-\arrlen,yshift=\y mm]sys.west) -- ([yshift=\y mm]sys.west);
}

% --- УПРАВЛЕНИЕ (сверху) ---
\foreach \x/\t in {
    -42/{Ограничения\\ предметной области},
    -14/{Требования\\ целостности данных},
    14/{Правила\\ разграничения доступа},
    42/{Математические модели\\ и алгоритмы расчёта\\ метрик}} {
  \draw[arr] ([xshift=\x mm,yshift=\stub]sys.north) -- ([xshift=\x mm]sys.north);
  \node[lbl, anchor=south, align=center, text width=24mm]
    at ([xshift=\x mm,yshift=\stub+1mm]sys.north) {\t};
}

% --- ВЫХОДЫ (справа) ---
\foreach \y/\t in {
    17/{Рассчитанные показатели\\ эффективности},
    6/{Рейтинги и лидерборды},
    -6/{Командные сводки и сравнение профилей},
    -17/{Турнирные таблицы}} {
  \draw[arr] ([yshift=\y mm]sys.east) -- ([xshift=\outlen,yshift=\y mm]sys.east);
  \node[lbl, anchor=west, align=left, text width=42mm]
    at ([xshift=\outlen+1mm,yshift=\y mm]sys.east) {\t};
}

% --- МЕХАНИЗМЫ (снизу) ---
\foreach \x/\t in {
    -34/{Веб-\\ приложение},
    -11.5/{База\\ данных},
    11.5/{Подсистема\\ кэширования},
    34/{Скрипт\\ загрузки\\ данных}} {
  \draw[arr] ([xshift=\x mm,yshift=-\stub]sys.south) -- ([xshift=\x mm]sys.south);
  \node[lbl, anchor=north, align=center, text width=22mm]
    at ([xshift=\x mm,yshift=-\stub-1mm]sys.south) {\t};
}
\end{tikzpicture}
\end{document}
